\documentclass[11pt]{article}
\usepackage[a4paper,margin=1in]{geometry}
\usepackage{fontspec}
\usepackage[boldfont=false]{xeCJK}
\setCJKmainfont{FandolSong-Regular.otf}
\usepackage{amsmath}
\usepackage{amssymb}
\usepackage{array}
\usepackage{booktabs}
\usepackage{graphicx}
\usepackage{adjustbox}
\usepackage{enumitem}
\setlist[itemize]{topsep=2pt,partopsep=0pt,parsep=0pt,itemsep=2pt,leftmargin=1.4em}
\setlist[enumerate]{topsep=2pt,partopsep=0pt,parsep=0pt,itemsep=2pt,leftmargin=1.6em}
\usepackage{parskip}
\usepackage{fvextra}
\fvset{breaklines=true,breakanywhere=true,fontsize=\small}
\setlength{\parindent}{0pt}

\begin{document}

\begin{center}
  {\LARGE\bfseries Organic synthesis}\\[0.35em]
  {\large A Level Chemistry}
\end{center}
\vspace{0.6em}

\section*{Organic synthesis}

This topic does not add new reactions. Instead it asks you to \textbf{join up} the reactions you already know, so you can build a target molecule in several steps.

\subsection*{Identifying functional groups}

A molecule may have more than one \textbf{functional group} 官能团. Use the test reactions from the syllabus to identify each one, and then predict how the molecule will behave. For example:

\begin{itemize}
\item decolourises bromine water → a C=C double bond (an \textbf{alkene} 烯烃).

\item gives a precipitate with silver nitrate → a \textbf{halogenoalkane} 卤代烷.

\item orange $\text{K}_2\text{Cr}_2\text{O}_7$ turns green → a primary or secondary \textbf{alcohol} 醇.

\item orange precipitate with 2,4-DNPH → an \textbf{aldehyde} 醛 or \textbf{ketone} 酮.

\item fizzes with a carbonate → a \textbf{carboxylic acid} 羧酸.

\end{itemize}

\subsection*{A map of the AS reactions}

Each row turns one functional group into another. Learn it as a map you can travel around:

\begin{center}
\adjustbox{max width=\linewidth}{%
\begin{tabular}{lll}
\toprule
\textbf{Start} & \textbf{Reagent and conditions} & \textbf{Product} \\
\midrule
alkene & $\text{H}_2$, Ni & alkane \\
alkene & $\text{HX}$, or $\text{X}_2$ & halogenoalkane \\
alkene & steam, $\text{H}_3\text{PO}_4$ & alcohol \\
halogenoalkane & $\text{NaOH(aq)}$, heat & alcohol \\
halogenoalkane & $\text{KCN}$ in ethanol, heat & a \textbf{nitrile} 腈 (adds one carbon) \\
halogenoalkane & $\text{NH}_3$ in ethanol, pressure & an \textbf{amine} 胺 \\
alcohol & $\text{K}_2\text{Cr}_2\text{O}_7$, distil / reflux & aldehyde / carboxylic acid \\
alcohol & concentrated acid, heat & alkene \\
aldehyde or ketone & $\text{NaBH}_4$ & alcohol \\
aldehyde or ketone & $\text{HCN}$, $\text{KCN}$ & hydroxynitrile \\
nitrile & dilute acid, heat & carboxylic acid \\
carboxylic acid + alcohol & concentrated $\text{H}_2\text{SO}_4$ & an \textbf{ester} 酯 \\
\bottomrule
\end{tabular}}
\end{center}

\subsection*{Planning a multi-step route}

To devise a \textbf{synthetic route} 合成路线:

\begin{enumerate}
\item compare the target with the starting material — what has changed (the functional group, the number of carbons)?

\item work \textbf{backwards} from the target: which single reaction could make it, and from what?

\item repeat until you reach the starting material.

\item write each step with its \textbf{reagent} 试剂 and conditions.

\end{enumerate}

If you need to add a carbon, the $\text{KCN}$ step is the key — it is the only AS reaction that lengthens the chain.

\subsection*{Analysing a route}

When you are given a route, for each step state the \textbf{type of reaction} (such as \textbf{oxidation} 氧化, \textbf{reduction} 还原, substitution, addition or elimination) and the reagent used. Also think about possible \textbf{by-products} 副产物 — for example, making an amine from a halogenoalkane also gives a mixture of further-substituted amines, so the yield of the simple amine is low.


\end{document}
